\chapter{Methodology}

\section{System overview}
The system has four layers: (1) log parser/normaliser, (2) feature
extractor, (3) rule engine, (4) ML model + hybrid decision fusion. Events
flow through the parser, then the rule engine and the ML model operate in
parallel before being fused into a final per-IP decision.

\section{Rule engine}
Re-implemented from the Capstone 1 abstract.
\begin{itemize}
  \item \textbf{Static rule:} $>\,N$ failed attempts from a single IP within
        a rolling $W$-second window.
  \item \textbf{Adaptive rule:} per-IP failure-rate baseline computed over the
        first $B$ days; alert whenever a one-minute bucket exceeds
        $\mu + k\sigma$.
  \item \textbf{Contextual rules:} invalid-user flood, user-spray
        (many unique users with few attempts each), low port diversity
        (script-driven attacks).
\end{itemize}

\section{Feature engineering}
For each IP the pipeline produces 14 features (counts, ratios, ratios of
failures to successes, Shannon entropy of usernames, port diversity,
burstiness, time-of-day, internal/external). The exact column list lives in
\texttt{src/features/build\_features.py}.

\section{Machine-learning models}
Random Forest, XGBoost (supervised) and Isolation Forest (unsupervised).
Class imbalance is mitigated using SMOTE from \texttt{imbalanced-learn}.

\section{Hybrid decision fusion}
Per-IP rule confidence and ML probability are combined as
$\alpha \cdot s_r + (1-\alpha) \cdot s_m$ where $\alpha \in [0, 1]$.
Either signal reaching its own block threshold immediately raises the
final decision to ``attack'' (fast path).

\section{Evaluation}
We compare static-only, adaptive-only, hybrid-rule-only, ML-only, and
hybrid (rule + ML) over the same dataset, reporting precision, recall,
F1, ROC-AUC, and per-event latency.
